\emph{Two tools.} (a) As in the proof of Theorem \claimref{synthesis.strict-completions-exist-exactly-when-the-signature-does-not-jump}: for a continuous family of self-adjoint maps, the spectral projection onto the eigenvalues lying in an interval whose end points are not eigenvalues is continuous, and its image is a vector bundle. (b) Let $\Pi(\sigma, s)$ be a continuous family of orthogonal projections in the fibres of a bundle over the compact space $\Sig$, $s \in [s_{0}, s_{1}]$. Then the bundles $\operatorname{im}\Pi(\cdot, s_{0})$ and $\operatorname{im}\Pi(\cdot, s_{1})$ over $\Sig$ are isomorphic. Indeed, by uniform continuity there is a subdivision such that $\|\Pi(\sigma,s) - \Pi(\sigma,s')\| < 1$ for $s, s'$ in one piece; then $\Pi(\sigma,s')$ is injective on $\operatorname{im}\Pi(\sigma,s)$ (a vector $v$ of that image killed by $\Pi(\sigma,s')$ has $|v| = |(\Pi(\sigma,s) - \Pi(\sigma,s'))v| < |v|$), the rank is constant in $s$ (the trace is continuous and an integer), so $\Pi(\cdot,s')$ restricts to a continuous isomorphism of the two images; compose over the pieces. This is the compactness argument of the proof of Theorem \claimref{synthesis.wall-class-one-rung-obstruction}. (c) The isomorphism classes of the eigenbundles of a nondegenerate form do not depend on the bundle metric: two metrics are joined by a segment of metrics, along which the eigenbundles vary continuously with constant rank, and (b), applied to the projections onto them that are orthogonal for the first metric, with the parameter of the segment in place of $s$, gives the isomorphism. So over the collar we use the pull-back of a metric of $V_{\Sig}$, and in (2) its sum with the pull-back of a metric of $W|_{\Sig}$.

\emph{(1).} Let $g(\sigma, s)$ be the self-adjoint map of $G$ in the fibre of $V_{\Sig}$ at $\sigma$. At $s = 0$ it has the eigenvalue $0$ with eigenline $\shL$, and its other eigenvalues are nonzero; the least of their absolute values is continuous and positive on the compact $\Sig$, so at least $2\varepsilon$ for some $\varepsilon > 0$. By uniform continuity there is $\delta > 0$ with $\|g(\sigma,s) - g(\sigma,0)\| < \tfrac12\varepsilon$ for $|s| \leq \delta$; eigenvalues of symmetric maps, in order, move by at most the norm of the change (Weyl's inequality, from the min-max characterisation; elementary, not in the store), so for $|s| \leq \delta$ the map $g(\sigma,s)$ has exactly one eigenvalue $\mu(\sigma,s)$ in $(-\varepsilon, \varepsilon)$, $p - 1$ eigenvalues above $\varepsilon$ and $q$ below $-\varepsilon$ (the counts at $s = 0$, where the restriction of $g$ to $\shL^{\perp}$ is the map of $\overline G$). By (a) the projections $\Pi_{0}, \Pi_{+}, \Pi_{-}$ onto these three parts are continuous on $\Sig \times [-\delta, \delta]$, with $\operatorname{im}\Pi_{0}|_{0} = \shL$ and $\operatorname{im}\Pi_{\pm}|_{0} = \overline G_{\pm}$. For $s \neq 0$ the form is nondegenerate, so $\mu \neq 0$; for $s > 0$ the signature $(p,q)$ forces $\mu > 0$ and for $s < 0$ the signature $(p-1,q+1)$ forces $\mu < 0$. Hence for $0 < s \leq \delta$, $G_{+}|_{s} = \operatorname{im}(\Pi_{+} + \Pi_{0})|_{s}$ and $G_{-}|_{s} = \operatorname{im}\Pi_{-}|_{s}$, and for $-\delta \leq s < 0$, $G_{+}|_{s} = \operatorname{im}\Pi_{+}|_{s}$ and $G_{-}|_{s} = \operatorname{im}(\Pi_{-} + \Pi_{0})|_{s}$. By (b) on $[0, \delta]$ and on $[-\delta, 0]$, $\operatorname{im}\Pi_{0}|_{\pm\delta} \cong \shL$ and $\operatorname{im}\Pi_{\pm}|_{\pm\delta} \cong \overline G_{\pm}$. On $\Sig \times [\delta, 1]$ and on $\Sig \times [-1, -\delta]$ the form is nondegenerate, the projections onto $G_{\pm}$ are continuous by (a), and (b) carries the statement from $|s| = \delta$ to every $s$ of the same sign.

\emph{(2).} Give $V \oplus W$ the sum of the metric of $V$ and a metric of $W$. By clause (4) of Theorem \claimref{synthesis.strict-completions-exist-exactly-when-the-signature-does-not-jump} (the transfer clause; its proof is (a) for the nondegenerate $h$), the eigenbundles $h_{\pm}$ are vector bundles over $M$, with continuous projections. By (b) on $[-1, 1]$, $h_{\pm}|_{1} \cong h_{\pm}|_{-1}$ over $\Sig$. At $s = \pm 1$ the form is an orthogonal sum for an orthogonal sum of metrics, so $h_{\pm}|_{1} = G_{\pm}|_{1} \oplus (K_{\mathrm{ph}})_{\pm}$ and $h_{\pm}|_{-1} = G_{\pm}|_{-1} \oplus (K_{\mathrm{far}})_{\pm}$. Insert (1).

\emph{(3).} The signature of $h$ is constant on the connected $M$, so $(p,q) + (k,k) = (p-1,q+1) + \operatorname{sig}K_{\mathrm{far}}$ and $\operatorname{sig}K_{\mathrm{far}} = (k+1,k-1)$. $H_{k}$ is a constant form on $\underline{\R}^{2k}$, so $(K_{\mathrm{ph}})_{\pm}$ are trivial of rank $k$, and the second isomorphism of (2) is the displayed one, with $N$ of rank $k - 1$. Taking top exterior powers, $\det\overline G_{-} \otimes \shL \otimes \det N \cong \det\overline G_{-}$, so $\shL \otimes \det N$ is trivial, and since $\shL \otimes \shL$ is trivial (\claimref{sheaf.imported-real-line-bundles}), $\det N \cong \shL$. For $k = 1$, $N = 0$ and $\shL$ is trivial. For $k = 2$, $N$ is a line, so $N = \det N \cong \shL$. By the Whitney formula (\claimref{synthesis.imported-whitney-product-formula-and-projective-space}), $w(\overline G_{-})\,(1 + x)\,w(N) = w(\overline G_{-})$; $w(\overline G_{-})$ has leading term $1$ and is invertible in the cohomology ring of the compact wall, so $(1 + x)\,w(N) = 1$, $w(N) = 1 + x + x^{2} + \cdots$, and since the classes of $N$ vanish above its rank $k - 1$ (the dimension axiom of Stiefel-Whitney classes, classical; the text of the imported Fact states the product formula and the triviality axiom and not this one, and the proof of the cited rung-count Theorem uses it in the same way), $x^{j} = 0$ for $j \geq k$. The comparison with Theorem \claimref{synthesis.wall-rung-count-is-height-of-dying-line-class}: its clause (1) asserts a bundle $N$ of rank $k - 1$ with an isomorphism of the same form, $N(\overline G)$ there being $\overline G_{-}$ here, for a form on $V \oplus \underline{\R}^{2k}$ over the physical side only that restricts to $G$ along the wall, where there is no far form; here the form is defined on both sides and nothing is asked along the wall. The two statements have different hypotheses and neither is deduced from the other.

\emph{(4).} $G = \pi_{\ell^{\perp}} + s\,\pi_{\ell}$ has the eigenvalue $s$ on $\ell$ and $1$ on $\ell^{\perp}$; its entries are quadratic in $(\cos\tfrac a2, \sin\tfrac a2)$, so they are continuous functions on the circle, while $\ell$ itself returns to $-\ell$ after one turn: $\shL = \R\ell$ is the twisted line, not trivial. (W) holds with $p = 2$, $q = 0$, $\overline G_{-} = 0$, $\overline G_{+} = \shL^{\perp}$. \emph{One rung.} The hyperbolic form is a constant form on the trivial plane, so its eigenbundles are trivial lines; the definite form on the other side may vary, its eigenbundle being the whole trivial plane. With the hyperbolic plane on the physical side, (3) for $k = 1$ would make $\shL$ trivial. With it on the far side, the physical form $K_{\mathrm{ph}}$ is negative definite on $\underline{\R}^{2}$ by the signature count, $(K_{\mathrm{far}})_{-}$ is a trivial line, and the second isomorphism of (2) reads $\underline{\R}^{2} \cong \shL \oplus \underline{\R}$; top exterior powers make $\shL$ trivial again. \emph{Two rungs.} The vectors $n$ and $b$ are orthonormal for every $t$ ($n \perp n^{\perp}$ in $\Pi$, $n \perp \ell$ since they lie in different summands, $Q(b) = \cos^{2}t + \sin^{2}t$), so $1 - 2(nn^{\mathsf T} + bb^{\mathsf T})$ is $-1$ on their span and $+1$ on its complement: signature $(4,2)$. Under $(\cos\tfrac a2, \sin\tfrac a2) \mapsto -(\cos\tfrac a2, \sin\tfrac a2)$ both $n$ and $b$ change sign and the form does not, so it is a continuous function of the point of the circle. At $t = 0$, $b = n^{\perp}$ and $nn^{\mathsf T} + bb^{\mathsf T} = \pi_{\Pi}$; at $t = \tfrac\pi2$, $b = \ell$. Let $t$ fall from $\tfrac\pi2$ to $0$ as $s$ rises from $-1$ to $1$. For $s \geq 1$ the form $((1-\lambda)G + \lambda) \oplus K_{\mathrm{ph}}$, $\lambda$ falling from $1$ to $0$, joins $1 \oplus K_{\mathrm{ph}}$ to $G \oplus K_{\mathrm{ph}}$ through forms of the same signature, $G$ being positive definite there; for $s \leq -1$, $((1-\lambda)(1 - 2\ell\ell^{\mathsf T}) + \lambda G) \oplus K_{\mathrm{far}}$ joins $(1 - 2\ell\ell^{\mathsf T}) \oplus K_{\mathrm{far}}$ to $G \oplus K_{\mathrm{far}}$, the eigenvalue on $\ell$ being $-(1-\lambda) + \lambda s < 0$. $K_{\mathrm{ph}} = 1 - 2\pi_{\Pi}$ is constant of signature $(2,2)$, the sum of two hyperbolic planes in a suitable constant basis; $K_{\mathrm{far}} = 1 - 2nn^{\mathsf T}$ has signature $(3,1)$ and negative line $\R n \cong \shL$. The form passes through the wall $s = 0$ without degenerating; there it does not restrict to $G$ on the plane bundle.

\emph{(5).} The class is the comparison of the slices $s > 0$ and $s < 0$ of the collar made in (1). From (1), the ranks: $\operatorname{rank}G_{+}$ falls by one and $\operatorname{rank}G_{-}$ rises by one, the signature jump, and $(p - q) - ((p-1) - (q+1)) = 2$. A real line bundle is determined up to isomorphism by $w_{1}$ and trivial exactly when $w_{1} = 0$ (\claimref{sheaf.imported-real-line-bundles}, with additivity of $w_{1}$ under tensor product: $\shL \otimes \shL'$ is trivial exactly when $w_{1}(\shL) = w_{1}(\shL')$, and then $\shL \cong \shL \otimes \shL' \otimes \shL' \cong \shL'$). If both exterior forms had one signature $(c,c')$, the first isomorphism of (2) would equate ranks $p + c$ and $p - 1 + c$: the stored obstruction.
